\documentclass[10pt,a4paper]{amsart}
\usepackage[margin=19mm]{geometry}
\usepackage{amsmath,amssymb,mathtools}
\usepackage[scr=rsfs,cal=euler]{mathalfa}
\usepackage{xfrac}
\usepackage[unicode,hidelinks,bookmarks=false]{hyperref}
\newcommand{\ie}{i.e.\ }
\newcommand{\cf}{cf.\ }
\newcommand{\resp}{respectively\ }
\newcommand{\Subsection}{\subsection}
\providecommand{\nobreakdash}{\nobreak\hbox{-}}
\setlength{\emergencystretch}{2em}
\multlinegap0pt
\usepackage{amssymb}
\usepackage{xcolor}
\usepackage[shortlabels]{enumitem}
\usepackage{tikz}
\newtheorem{lemma}{Lemma}
\newcommand{\Z}{\mathbb{Z}}
\title{The Profinite Rigidity of Torsion-Free Lamplighter Groups}
\author{Nikolay Nikolov, Julian Wykowski}
\date{Source: arXiv:2512.19295v1 (CC BY 4.0)}
\begin{document}
\maketitle

\noindent\emph{Source and licence.} The Profinite Rigidity of Torsion-Free Lamplighter Groups. Version arXiv:2512.19295v1 (CC BY 4.0), distributed by arXiv under the Creative Commons Attribution 4.0 International licence, http://creativecommons.org/licenses/by/4.0/, as recorded in the arXiv licence metadata for this exact version and shown on the abstract page https://arxiv.org/abs/2512.19295v1. Full licence notice: \texttt{licenses/arxiv-2512.19295-CC-BY-4.0.txt}.\par\smallskip

\noindent\textit{Editorial scope.} Counted unit: the article's lemma Lem::Abelian (source0.tex lines 295--296). The article's complete original proof of it is delivered with it (source0.tex lines 298--315, one \texttt{proof} environment of 18 lines). No mathematics is written, repaired or re-derived: the statement and the proof are the article's own lines, in the article's own notation, and no retained line is substituted or re-flowed.  Retained uncounted as the source-local context the proof needs: lines 88--90; lines 265--266.  Nothing was omitted from any retained span, so the package carries no omission record.  No retained line contains a citation, so the wrapper carries no bibliography and no bibliography entry.  No retained line contains a figure environment, an image-inclusion command or drawing code, so no image is delivered and the package declares no asset dependency.  Editorial formatting only, in addition: an AMS \texttt{amsart} preamble that restates the article's own declarations and macros used by the retained lines, namely the environments \texttt{lemma} on separate counters, together with the standard packages those lines need.  The retained environments print in this document as Lemma 1 in order of appearance, whereas the article prints the same items with the numbering of its own full document.  The complete native source of the article as distributed in the arXiv e-print for this exact version is bundled unchanged as \texttt{sources/191459/source0.tex}.
\begin{equation} \label{Eq::DefOfLL}
\Gamma_n = \Z^n \wr \Z \cong M \rtimes \langle t\rangle
\end{equation}

\begin{lemma} \label{centralizer} The subgroup $N \trianglelefteq \Delta$ centralises the derived subgroup $\gamma_2(\Delta)$ of $\Delta$.
\end{lemma}

\begin{lemma}\label{Lem::Abelian} The group $N$ is abelian.
\end{lemma}

\begin{proof} We commence by exhibiting a cofinal family of finite quotients of the torsion-free lamplighter group $\Gamma = \Z^n \wr \Z$. Consider the semidirect product decomposition of $\Gamma$ in (\ref{Eq::DefOfLL}) and fix generators $a_1, \ldots, a_n$ for the free $\Z[\langle t \rangle]$-module $M$, so that $\Gamma$ is as a group by $a_1, \ldots, a_n$ and $t$. Given a positive integer $m \in \mathbb{N}$, define
\[
\Gamma(m) = C_m^n \wr C_m \cong R_m^n \rtimes C_m
\] where $R_m$ is the additive abelian group of the ring
\[
R_m = \frac{(\Z/m\Z)[x]}{(x^m - 1)}
\]
and the generator $\tau_m$ of $C_m$ acts on each copy of $R_m$ as multiplication by $x$. The natural maps $M \twoheadrightarrow R_m$ and $\langle t \rangle \twoheadrightarrow C_m$ yield an epimorphism $\pi_m \colon \Gamma \twoheadrightarrow C_m$. In fact, given any finite-index normal subgroup $U \trianglelefteq \Gamma$, setting $m = 2\cdot [\Gamma : U]$ and $V = \langle \langle t^m,a_1^m,\ldots, a_n^m\rangle \rangle \trianglelefteq \Gamma$ yields $V \leq U$ and $\Gamma / V \cong \Gamma(m)$. Hence the set $\{\Gamma(m) : m \in \mathbb{N}\}$ is a cofinal subsystem of the inverse system of finite quotients $\mathcal{C}(\Gamma) = \mathcal{C}(\Delta)$. To prove that $N$ is abelian, it thus suffices to show that $\pi_m(N) \leq \Gamma(m)$ is abelian for $m > 2$.
Moreover, Lemma~\ref{centralizer} yields $\pi_m(N) \leq C_{\Gamma(m)}([\Gamma(m), \Gamma(m)])$, so in fact it will suffice to show that the centraliser $Z(m) = C_{\Gamma(m)}([\Gamma(m), \Gamma(m)]) \trianglelefteq \Gamma(m)$ is abelian for each positive integer $m > 2$.

Indeed, we shall demonstrate that $Z(m) = R_m^n$, which is abelian by construction. On the one hand, we have $[\Gamma(m), \Gamma(m)] = (x-1) R_m^n$ which is centralised by its abelian supergroup $R_m^n$, so $R_m^n \leq Z(m)$. Conversely, suppose for a contradiction that there exists $z \in Z(m) \backslash R_m^n$. We may assume without loss of generality that $z=\tau^k$ for some $k \in \{1, \ldots, m-1\}$. The condition that $z$ centralizes $[\Gamma(m),\Gamma(m)]$ is equivalent to $\tau^k-1$ annihilating $(x-1)R_m^n$, which in turn is equivalent to the equation $(x^k-1)(x-1)=0$ in $R_m$. Now, if $k \neq 1$ then
\[
(x^k-1)(x-1)=x^{k+1}-x^k-x+1 \neq 0
\] because $x \not \in \{1,x^k,x^{k+1}\}$ and $1,x, \ldots x^{m-1}$ are a basis of $R_m$ as a free $\Z/m\Z$-module. On the other hand, if $k = 1$ then \[
(x^k-1)(x-1)=x^2 -2x+1 \not =0
\]
because $m>2$. This contradicts $z \in Z(m)$, so we conclude that $Z(m) = R_m^n$ is abelian and the proof is complete.
\end{proof}

\end{document}
